\documentclass{standalone}
\usepackage{tikz}
\usepackage{pgfplots}
\usepgfplotslibrary{groupplots, fillbetween}
\pgfplotsset{compat=1.17}
\usepackage{times}
\usepackage[utf8]{vietnam}

% Hàm gauss(mu, sigma)
\pgfmathdeclarefunction{gauss}{2}{%
	\pgfmathparse{1/(#2 * sqrt(2 * pi)) * exp(-((x - #1)^2) / (2 * #2^2))}%
}

\begin{document}
	
	\begin{tikzpicture}
		\begin{groupplot}[
	group style={group size=3 by 1, horizontal sep=1cm},
	width=6cm, height=6cm,
	domain=-5:10, samples=100,
	anchor=south west,
	title style={at={(0.5,-0.15)}, anchor=north}, % đặt tiêu đề phía dưới
	]

			
			% ===== Frame 1: min(f,g) =====
			\nextgroupplot[title={(a)}]
			\addplot [color=black!45, line width=1.2pt, dashed] {gauss(1,1)};
			
			\addplot[color=black!45, line width=1.2pt, dashdotted]{gauss(4,1)};
			
			\addplot [color=black!95, line width=2pt, fill=black!45, opacity=.5]  {min(gauss(1,1), gauss(4,1))};
			
			% ===== Frame 2: (f-g)^2 =====
			\nextgroupplot[title={(b)}]
			\addplot [color=black!45, line width=1.2pt, dashed] {gauss(1,1)};
			
			\addplot [color=black!45, line width=1.2pt, dashdotted]{gauss(4,1)};
			
			\addplot [color=black!95, line width=2pt, fill=black!45, opacity=.5] {(gauss(1,1) - gauss(4,1))^2};
			
			% ===== Frame 3: sqrt(f g) =====
			\nextgroupplot[title={(c)}]
			\addplot  [color=black!45, line width=1.2pt, dashed] {gauss(1,1)};
			
			\addplot [color=black!45, line width=1.2pt, dashdotted] {gauss(4,1)};
			
			\addplot  [color=black!95, line width=2pt, fill=black!45, opacity=.5]  {sqrt(gauss(1,1) * gauss(4,1))};
			
		\end{groupplot}
	\end{tikzpicture}
	
\end{document}
